% G04 (issue #60): Section 2, outline of the argument.
% Author draft pass 1, session ADB991-claude-20261113-15a6; not reviewed.
% Not \input anywhere: placement after the statement section is reserved
% for F01 (see blueprint/units/YM-000026/handoff/results/G04/).
\section{Outline of the argument}
\label{ex-outline}

This section describes the route from the Wilson integrals of
Section~\ref{ex-lattice} to the target of Section~\ref{ex-statement}.
It proves nothing: each claim below is a pointer to the section where it is
proved or where it is stated as an outstanding input. The order of the
discussion follows the logical dependence of the steps, which differs in
places from the order of presentation. Reading the outline does not replace
reading the hypotheses of the cited results.

\subsection{Status of the argument}
\label{ex-outline-status}

The final step, Proposition~\ref{ex-conditional-existence}, deduces the
target from one construction statement, Input~\ref{ex-joint-input}.
Everything else in the article serves to establish that input. It requires
four things simultaneously for one lattice family and one extracted
hierarchy: the Wilson limits along the full mesh sequence, the phase
decomposition, the reconstructed relativistic theory, and the
operator-product and trace properties. At the present draft several
producers of this input remain conditional. The precise open interfaces are
listed in Section~\ref{ex-outline-open}. In particular, this article does not
yet contain a complete existence proof, and nothing in the present section
should be read as asserting one.

Two scope limitations of the target itself should be kept in view. The
target concerns the constructed gauge-invariant fields; it does not
identify a given label with a separately prescribed physical observable.
It contains no mass-gap statement. The short-distance trace test of
Section~\ref{ex-assembly} excludes finite combinations of mass-shell
profiles for one two-point function. That is a non-triviality statement,
not a spectral-gap statement.

\subsection{Prerequisites}
\label{ex-outline-prerequisites}

Section~\ref{ex-statement} lists the prerequisites assumed of the reader.
The detailed proofs also use the following standard material without
proof: Fubini--Tonelli, Dynkin's $\pi$--$\lambda$ theorem and Bochner
integrals of Banach-space-valued functions; Cauchy's theorem and integral
formula; the Hahn--Banach theorem; the Fourier transform on Schwartz space;
and the elementary inequalities for Gaussian vectors (Markov, Chernoff and
Jensen inequalities, Wick's rule). Commutative $C^*$-algebra theory, the
Riesz--Markov theorem, the spectral theorem for unitary representations of
$\mathbb R^4$, the topology of the Schwartz space and the local structure of
closed matrix groups are cited at their points of use. They enter only in
Sections~\ref{ex-reflection}--\ref{sec:lorentz-covariance} and in the
nearest-point projection onto $G$ (Lemma~\ref{lem:branch-nearest-projection}).

\subsection{Context}
\label{ex-outline-context}

The finite-volume measures are Wilson's lattice gauge theory
\cite{Wilson74}. The bare coefficients grow logarithmically,
\eqref{ex-running-input}, as asymptotic freedom predicts
\cite{GrossWilczek73,Politzer73}. The construction proceeds by block-spin
renormalization for lattice gauge fields, an approach developed in
Ba\l{}aban's programme \cite{Balaban87RG1}. The present argument follows the
pinned formal proof rather than those papers. No theorem of
\cite{Balaban87RG1} is used, and where a construction resembles one there,
the resemblance is noted without being relied upon. The passage from
Euclidean correlations to a relativistic theory has the shape of the
Osterwalder--Schrader reconstruction \cite{OS73,OS75}. Here it is carried
out directly, for the specific hierarchy and with its own estimates, in
Sections~\ref{ex-reflection}--\ref{sec:lorentz-covariance}. The comparison
with the corrected theorem of \cite{OS75} remains, as stated in
Section~\ref{ex-reflection}, one possible literature route, not an input.

\subsection{The order of the limits}
\label{ex-outline-limits}

The target is stated in Section~\ref{ex-statement}. Its essential feature is
an order of quantifiers. One regular lattice family is chosen first, with
spacings $a_N=D_N^{-1}$, periods $P_p$ and bare coefficients
$\beta_{N,p}$ satisfying \eqref{ex-running-input}. Only afterwards is a
volume sequence $p(k)$ prescribed. For it a subsequence $q(k)$ and a
hierarchy of graded fields, phases and correlations are chosen. The
ultraviolet limit \eqref{ex-row-target} is taken along the \emph{full} mesh
sequence at each fixed period. The infinite-volume limit
\eqref{ex-volume-target} is taken along the subsequence. The same
subsequence and hierarchy must serve every arity, token and test, and also
the operator-product and trace statements.

Lemma~\ref{ex-full-mesh} shows how a full mesh limit is combined with a
hierarchy obtained by compactness. It does not produce full mesh
convergence from subsequential convergence. That must be proved separately,
and is the content of Section~\ref{sec:fullmesh-correlation-cauchy}.

\subsection{The plan in nine steps}
\label{ex-outline-plan}

The argument has nine main steps. Steps~1--4 construct the lattice
estimates. Steps~5--6 take the limits and obtain the Euclidean theory and
its phases. Step~7 reconstructs the relativistic theory. Steps~8--9 establish
the operator-product expansion, the stress tensor and the trace test. The
final assembly is Proposition~\ref{ex-conditional-existence}. The table in
Section~\ref{ex-outline-map} maps the steps to sections.

\paragraph{Step 1: the free block-spin transformation.}
The Wilson measure on a periodic lattice is gauge invariant and reflection
positive for both site and link reflections
(Propositions~\ref{prop:ex-lattice-gauge}
and~\ref{prop:ex-lattice-finite-reflection}). Together with the lattice
symmetries, reflection positivity is what passes to the limit and yields
the Euclidean Hilbert space of Step~6.

The block-spin renormalization group averages over blocks of odd side $L$
and changes the spacing by a factor $L$ per step. Its quadratic part is
analysed first, uniformly in the period. The fine-to-coarse transports and
the free precisions of each scale are given by explicit Laurent symbols.
Contour shifts give exponential decay of their kernels, uniform in the
period (Theorem~\ref{thm:ex-lattice-wire-tower-strip} and
Proposition~\ref{prop:ex-lattice-wire-entrywise-bound}).
Section~\ref{ex-lattice} also establishes the exact algebraic identities of
the transports on every torus, including zero momentum.

The quadratic forms at different starting scales differ by a neutral,
translation-covariant representative. One block step contracts this
difference by a factor $L^{-1}$, up to explicit constants
(Proposition~\ref{prop:ex-lattice-flat-step}). An elementary recurrence,
Proposition~\ref{prop:ex-lattice-flat-rate}, converts the one-step
contraction into power decay over a finite interval of scales, uniformly
in the starting level (Theorem~\ref{thm:ex-lattice-rg-construction}).

\paragraph{Step 2: one nonlinear block step.}
A block step of the interacting measure is analysed on a fixed sector, that
is, for a fixed choice of the regions where the field is large. Its
analysis has four parts.
\begin{enumerate}
\item \emph{Defects and their energy.} Away from the large-field region,
 gauge comparison through axial gauges and minimising frames bounds the
 curvature of the coarse predictor. The predictor uses the analytic
 nearest-point projection onto $G$
 (Lemma~\ref{lem:branch-nearest-projection}). Each defect seed is then
 assigned a definite share of action, which pays for it.
\item \emph{The Gaussian part.} The precision of the selected rows is
 uniformly positive on each sector, with a separate local gauge at each
 active bond. Its inverse is exponentially localized and approximated by
 principal-window inverses, with explicit control of the log-determinant
 (Section~\ref{sec:local-covariance}). The Haar density in exponential
 coordinates, a change of variables holding the fine field fixed, and the
 jets of the local actions are controlled. Together they give a
 differentiated bound for the product of local activities.
\item \emph{Exact reassembly.} The block integral is rewritten exactly as a
 completed-square Gaussian with one scalar overhead, multiplied by a gas of
 compatible components (Section~\ref{sec:gaussian-reassembly}). Averaging
 over the forests separates the regular part. The exclusion correction and
 the localized logarithms are absorbed into a gas of forced activities.
 A marked insertion is carried through the step exactly, as an affine tilt.
 Its derivative yields the elementary projection estimate used by the
 operator-product construction (Section~\ref{sec:projection-payment-producer}).
\item \emph{Cluster expansion.} A tree-graph bound for repeated-label
 cluster expansions gives the logarithm of the polymer gas as a convergent
 sum of connected terms, without any positivity assumption on the
 activities (Section~\ref{sec:cluster-expansion}). Warning events of the
 Gaussian fields are controlled by a union bound
 (Proposition~\ref{prop:warning-crowding}).
\end{enumerate}
The scale-transfer of local representations
(Section~\ref{sec:local-transfer}) and exact jet extraction
(Section~\ref{sec:local-splitting}) complete the one-step map.

\paragraph{Step 3: stationary analysis and the one-loop coefficient.}
On the small-field region the block integral is evaluated around a
stationary configuration. A contraction argument gives a stationary branch
on a ball whose radius and Lipschitz constant do not depend on the volume
(Section~\ref{sec:uniform-stationary-chart}). The truncated Gaussian
integral about it has a logarithm of order $\beta^{-1}$
(Section~\ref{sec:truncated-stationary-log}). The remainder of the pure
nonlinear push is smaller by a definite power
(Section~\ref{sec:pure-action-small-field}).

The quadratic coefficient of the output action is then computed. The
cubic and quartic classical vertices come from stationary elimination
(Proposition~\ref{prop:ex-classical-stationary-pictures}). The one-loop
determinant is compared with a reference determinant
(Section~\ref{sec:ex-thermal-contact-cancellation}). An annulus estimate
isolates the logarithmic term. The coefficient of $\log L$ in the change of
the inverse coupling over one step is
\[
 \lambda_G=\frac{11C_G}{24\pi^2},
\]
up to a constant independent of $L$
(Proposition~\ref{prop:ex-annulus-estimate}). Here $C_G$ is the adjoint
Casimir. This is the step at which the one-loop coefficient of asymptotic
freedom enters the construction. The canonical running-coefficient unit
writes the same coefficient as $2b_0(G)$ with
$b_0(G)=11C_{\mathfrak g}/(48\pi^2)$, and proves $b_0(G)>0$ from the
noncommutativity of $\mathfrak g$. That unit is not yet included in this
article, and checking that the two Casimir normalizations agree is part of
the running interface listed in Section~\ref{ex-outline-open}.

\paragraph{Step 4: the coupling history and the common funded history.}
The one-step outputs are attached to paths of quadratic backgrounds. Their
limits as the cutoff is removed at a fixed scale are unique
(Section~\ref{sec:background-path-convergence}). A real-coupling
interpolation gives analytic parameter estimates
(Section~\ref{sec:real-coupling-interpolation}). The coupled coupling and
activity histories then converge (Section~\ref{sec:coupled-background-convergence}).

The coupling is reset at each step by a scalar shooting argument through
bands of admissible values. This keeps the inverse coupling within a
sublinear error of a line in $s_m=\log a_m^{-1}$
(Section~\ref{sec:surviving-interval-shooting}). Its slope is a
group constant $s_G>0$. Once $s_G$ is matched with the one-loop coefficient,
the running-coefficient unit converts this estimate into the uniform
logarithmic growth \eqref{ex-running-input} of the public family. The terminal step installs the exact Wilson density
(Proposition~\ref{prop:wilson-startup-family}).

All finitely many numerical requirements at each stage are then paid on one
scale schedule. The block sides are chosen first, then their costs, then a
slowly advancing stage clock (Section~\ref{sec:common-scale-selection}).
This produces a single \emph{funded history}: one sequence of scales,
couplings and local representations on which the stationary, nonlinear and
observable estimates hold together (Section~\ref{sec:common-funded-history}).
The point of this step is that later limits use the estimates of one fixed
history. They do not use separate histories selected for each estimate.

\paragraph{Step 5: centered bounds and the full mesh limit.}
Local observables are constructed on the funded history, together with
their calibration. The microscopic equations of motion come from Haar
integration by parts on one link (Section~\ref{sec:marked-campaign-producer}).

The basic estimate is a root-test bound for a local observable copied into
every cell of a tiling. Reflection positivity and the chessboard inequality
turn it into control of all separated correlations
(Theorem~\ref{ex-centered-producers}). On the funded history the root
test is bounded by $Aa_m^{-d-\varepsilon}$ for a field of grade $d$,
uniformly in the period, the site and the label in a bounded set
(Theorem~\ref{ex-centered-main-bound}). This yields equicontinuity and
local bounds for all separated correlations.

The correlations converge along the full mesh sequence at each fixed period
(Section~\ref{sec:fullmesh-correlation-cauchy}). This uses a fixed
intermediate level at which the higher marked terms are uniformly small,
together with convergence at a fixed scale
(Section~\ref{sec:fixed-rung-convergence}). Bounds on lattice difference
quotients make the limits smooth on separated configurations
(Section~\ref{ex-smooth-limits} and
Theorem~\ref{thm:ex-fullmesh-fixedbounds-producer}). A diagonal extraction
over a countable basis then gives the infinite-volume table
(Proposition~\ref{prop:ex-limits-extraction}). Symmetry, Euclidean
covariance under the lattice symmetries, and reflection positivity pass to
the limit (Proposition~\ref{prop:ex-limits-euclidean}).

\paragraph{Step 6: the Euclidean Hilbert space and the phases.}
Reflection positivity gives a Hilbert space with a self-adjoint contraction
semigroup in Euclidean time and commuting spatial translations
(Theorem~\ref{ex-euclidean-semigroup}). The infinite-volume limit need
not have a unique vacuum. The time-invariant subspace carries a commutative
algebra generated by compressed slab insertions. Its boundedness rests on
uniform bounds for arbitrarily many slab copies
(Section~\ref{sec:actual-slab-copy-bounds}).

The Gelfand and Riesz--Markov theorems \cite{Williams,Sawyer} give a
probability measure $\nu$ of full support on the characters of this
algebra. The mixed correlations are the $\nu$-average of the character
correlations; this is the mixture identity~\eqref{ex-mixture}
(Proposition~\ref{ex-character-measure}). For $\nu$-almost every character
the ground space is one-dimensional (Theorem~\ref{ex-pure-ground}). The
selected character $\chi_*$ is chosen among these. Every character keeps the
regularity bounds, the null relations of the original fields
(Lemma~\ref{ex-phase-original-null}) and full $O(4)$ covariance
(Proposition~\ref{ex-individual-rotations}). The field spaces are therefore
the same in every phase (Proposition~\ref{ex-no-new-null}).

\paragraph{Step 7: analytic continuation and reconstruction.}
Functional calculus continues the heat semigroup to complex time
(Proposition~\ref{prop:complex-heat}) \cite{Williams}. Ordered insertions of
centered slab operators give holomorphic functions on the ordered tube.
Their tempered boundary values are the Wightman distributions
(Theorem~\ref{thm:joint-boundary}). Exponential testing places their
Fourier support in a half-space (Proposition~\ref{prop:half-support}).

Calibrated finite-dimensional Lorentz actions on the field levels descend
to the field union (Input~\ref{input:calibrated-actions} and
Proposition~\ref{prop:lorentz-label-descent}). A boost Ward identity then
gives Lorentz covariance (Theorem~\ref{thm:lorentz-covariance}). This
improves the support to the forward cones and gives holomorphy on the full
tube (Proposition~\ref{prop:full-tube-holomorphic}).

The fields are built on the span of words applied to the vacuum, with the
unitary Poincar\'e representation of Proposition~\ref{prop:word-motion}.
Nonzero labels act nontrivially (Proposition~\ref{prop:boundary-detection})
\cite{Melrose}. The joint spectral measure of the translations lies in
$\overline V_+$, and the translation-invariant vectors are the multiples of
the vacuum (Theorem~\ref{thm:field-spectrum}) \cite{BHV}. Spacelike
commutativity follows by continuing along imaginary boosts into both
ordered tubes (Theorem~\ref{thm:spacelike-exchange}). The boundary-value
identity~\eqref{eq:public-selected-boundary} relates the selected Euclidean
correlations to the vacuum expectation values of the reconstructed fields.

\paragraph{Step 8: the operator-product expansion.}
The limiting fields carry a grading, and the operator-product statement
\eqref{ex-ope-coefficient}--\eqref{ex-ope-remainder} is uniform over all
phases. One prescription of fields and finite exterior tests is fixed for
every truncation degree (Proposition~\ref{prop:ope-prescription}). Its raw
coefficients are averaged over $O(4)$ and split into homogeneous parts.

Bounds pass to every phase through an even-replica chessboard argument and
the full support of $\nu$ (Section~\ref{sec:ope-phase-prices}). Lattice
clusters are grouped exactly in a square-free mark algebra. The grouped
local field at the last scale is matched with a finite-cutoff reference
observable by a backward cone argument (Section~\ref{sec:reference-match}).
Under the local matching estimates (Input~\ref{input:ope-match}, assembled
in Proposition~\ref{prop:ope-match-assembly}) and the raw product
comparison (Input~\ref{input:ope-raw-estimates}),
Theorem~\ref{thm:ope-quantitative} gives the coefficient and remainder
bounds. Their thresholds are uniform in the phase.

\paragraph{Step 9: the stress tensor and the trace test.}
The field quotient of the interacting theory is identified with that of the
free Maxwell theory through finite rank. Finite point probes detect each
grade in the free theory (Theorem~\ref{thm:ex-lattice-all-grade-free-detector}).
Comparisons purchased on the funded history transport the kernels of the
Maxwell equations (Section~\ref{sec:finite-rank-kernel-transport}). The
supporting comparisons of interacting and Gaussian ``reader'' fields over
short intervals of scales are in Sections~\ref{sec:restarted-rank-history}--%
\ref{sec:selected-history-limits}.

This gives one conserved physical stress tensor with a fixed nonzero
normalization along one comparison curve
(Section~\ref{sec:actual-stress-normalization}). Its translation and rotation
Ward identities hold at all separated configurations
(Theorem~\ref{thm:ward-global-propagation}). They hold in every phase
(Proposition~\ref{prop:phase-ward-inheritance}). The same comparison
supplies the calibrated Lorentz actions used in Step~7
(Section~\ref{sec:physical-angular-actions}).

The centered trace $\Theta$ of the stress tensor has a positive Laplace
representation in the selected phase (Section~\ref{sec:trace-spectrum}).
Its short-distance asymptotic~\eqref{ex-trace-limit} is reduced, in
Theorem~\ref{thm:trace-asymptotic}, to the trace comparison along scales
(Input~\ref{input:trace-scales}). There the amplitude
$3\dim_{\mathbb R}\mathfrak g/\pi^4$ comes from the free contraction of the
curvature-square field and the Ward normalization. The factor
$\log^{-2}(r^{-1})$ comes from the running coupling with coefficient $b_0$.
Theorem~\ref{thm:profile-exclusion} then shows, without further input, that
such an asymptotic excludes finite combinations of mass-shell profiles.

\subsection{Map of the article}
\label{ex-outline-map}

The table gives, for each step, the sections where it is carried out, in
order of presentation. The soft analytic layer (Steps~5--8) is presented
before the construction estimates (Steps~2--4 and~9) that supply its
quantitative inputs.

\begin{center}
\small
\begin{tabular}{p{0.06\linewidth}p{0.36\linewidth}p{0.48\linewidth}}
\hline
Step & Content & From section \(\to\) to section\\
\hline
1 & lattice, free RG, finite-interval contraction &
 \ref{ex-lattice} \(\to\) \ref{sec:finite-menu-history-handoff}\\
5 & centered bounds; full mesh and volume limits &
 \ref{ex-centered-estimates} \(\to\) \ref{ex-public-same-run-cauchy}\\
6 & reflection, Hilbert space, phases &
 \ref{ex-reflection} \(\to\) \ref{ex-phase-fields}\\
7 & continuation, covariance, reconstruction &
 \ref{sec:complex-heat} \(\to\) \ref{sec:locality}\\
8 & OPE prescription, grouping &
 \ref{sec:ope-prescription} \(\to\) \ref{sec:ope-singletons}\\
2 & one nonlinear block step &
 \ref{sec:local-transfer} \(\to\) \ref{sec:high-degree-cluster-prices}\\
3 & stationary analysis, one-loop coefficient &
 \ref{sec:local-parameter-comparison} \(\to\) \ref{sec:ex-quadratic-pole-transport}\\
4 & paths, coupling reset, funded history &
 \ref{sec:free-horizon-convergence} \(\to\) \ref{sec:actual-phase-grid-approximants}\\
9 & reader comparisons and rank limits &
 \ref{sec:restarted-rank-history} \(\to\) \ref{sec:selected-history-limits}\\
8 & OPE matching and estimates &
 \ref{sec:ope-local-projection} \(\to\) \ref{sec:ope-estimates}\\
9 & free detectors, stress, Ward identities, trace &
 \ref{sec:ex-lattice-trace-free-rank} \(\to\) \ref{sec:ex-lattice-trace-local-charge}\\
4 & schedule purchases and binding &
 \ref{sec:nonlinear-payment-selection} \(\to\) \ref{sec:public-run-binding}\\
--- & assembly & \ref{ex-assembly}\\
\hline
\end{tabular}
\end{center}

\subsection{Open interfaces}
\label{ex-outline-open}

At this draft Input~\ref{ex-joint-input} is not yet established. The
following interfaces between the steps remain to be closed. Each is stated
precisely at its point of use.
\begin{enumerate}
\item \emph{Steps 2--4.} The one-step theorem is not yet assembled from its
 parts, so the stationary paths are attached only conditionally. The
 annulus estimate depends on three stated response and ticket premises
 (Proposition~\ref{prop:ex-annulus-estimate}). The slope $s_G$ of the
 running line in Step~4 must be matched with $\lambda_G$ and with the
 hypotheses of the running-coefficient unit, including its Casimir
 normalization.
\item \emph{Step 4 to Steps 5 and 9.} The selected public run must be bound
 to the funded history used by the centered service and the local course.
 The served representation used by the marked calibration must be produced.
\item \emph{Step 5.} Full mesh convergence must be matched, on the same
 history, with the subsequence refinement used for the fields.
\item \emph{Steps 6--7.} The field levels and their inclusions are used but
 not yet constructed by a named producer. The calibrated actions of
 Input~\ref{input:calibrated-actions} require the stress comparison of
 Step~9. The selected phase must be one of the $\nu$-almost every
 characters with a unique ground state.
\item \emph{Step 8.} Inputs~\ref{input:ope-match}
 and~\ref{input:ope-raw-estimates} require the matched fields of
 Steps~2, 4 and~5 and the physical filtration of Step~9 on the same labels.
\item \emph{Step 9.} The trace comparison of Input~\ref{input:trace-scales},
 the weighted sensor estimates and the branch-moment limits are supplied
 outside the included chapters and are still to be integrated.
\end{enumerate}
A final independent check of the union of these interfaces, for
circularity and for agreement of common data, is required before the
existence statement can be asserted.
